% !TITLE 桃花庵歌
% !AUTHOR 唐寅
% !DYNASTY 明清
% !GENRE 七言古诗
% !ORDER 1

\begin{poem}
桃花坞\textsuperscript{1}里桃花庵\textsuperscript{2}，桃花庵下桃花仙。
桃花仙人种桃树，又摘桃花换酒钱。
酒醒只在花前坐，酒醉还来花下眠。
半醉半醒日复日，花落花开年复年。
但愿老死花酒间，不愿鞠躬车马前\textsuperscript{3}。
车尘马足\textsuperscript{4}富者趣\textsuperscript{5}，酒盏花枝贫者缘\textsuperscript{6}。
若将富贵比贫者，一在平地一在天。
若将贫贱比车马，他得驱驰我得闲。
别人笑我忒\textsuperscript{7}风颠\textsuperscript{8}，我笑他人看不穿\textsuperscript{9}。
不见五陵豪杰墓\textsuperscript{10}，无花无酒锄作田。
\end{poem}

\begin{annotations}
\notetext{1}{桃花坞：苏州城北的一处地名，唐寅晚年在此地建桃花庵居住。坞（wù），四面高中间低的地方。}
\notetext{2}{庵：小茅屋、草舍。不是尼姑庵的庵，指隐士居住的简陋房屋。}
\notetext{3}{鞠躬车马前：在达官贵人的车马前弯腰行礼。鞠躬，弯腰致敬。}
\notetext{4}{车尘马足：车马的尘土和马蹄的奔忙，代指富贵人家的应酬奔波。}
\notetext{5}{趣：通“趋”，追逐、奔赴。富人们追逐着车马红尘中的名利。}
\notetext{6}{缘：缘分、命分。贫者与花酒为伴，也是一种命里注定的生活方式。}
\notetext{7}{忒：太、过于。忒（tuī），程度副词。}
\notetext{8}{风颠：疯狂、癫狂。“风”通“疯”。唐寅自称行为放诞不合世俗。}
\notetext{9}{看不穿：看不透。指世人看不穿名利富贵的虚妄。}
\notetext{10}{五陵豪杰墓：五陵是汉代五座帝王陵墓所在地（长陵、安陵、阳陵、茂陵、平陵），附近多住豪门贵族。那些豪杰的坟墓，如今变成了种庄稼的田地。}
\end{annotations}

\begin{appreciation}
唐寅是明代的苏州才子，乡试第一，人称唐解元，本来前程无量，后来卷进一桩科场舞弊案，一辈子和官场无缘。晚年他在苏州城外的桃花坞盖了间桃花庵，种桃花，换酒喝，写下这首《桃花庵歌》——中国诗里最潇洒的自我介绍之一。

开头像绕口令："桃花坞里桃花庵，桃花庵下桃花仙。"桃花一个词接一个词，句子像民歌一样朗朗上口。他是故意的——用土气热闹的调子开场，好像在说：我这个过闲日子的人，用不着典雅的腔调。"桃花仙人种桃树，又摘桃花换酒钱"——桃花哪能卖钱，真正换钱的是他的画，他偏说桃花，把谋生说成风雅。

"酒醒只在花前坐，酒醉还来花下眠。半醉半醒日复日，花落花开年复年。"人过一天，花开一年，都在打转，也都圆满。接着是选择："但愿老死花酒间，不愿鞠躬车马前。"要么醉死在花下，要么弯着腰站在贵人的车马前——他选前者。他有资格说这句话：他本来可以站在车马前，被人踢出来之后，反而看清了。"若将富贵比贫者，一在平地一在天"——仿佛富贵在天上；可马上补一句"若将贫贱比车马，他得驱驰我得闲"——天上的人在跑，地上的人在歇，谁舒坦还真不好说。

"别人笑我忒风颠，我笑他人看不穿。"世人笑他疯，他笑世人看不穿——看不穿名利的虚妄。末句是杀手锏："不见五陵豪杰墓，无花无酒锄作田。"汉朝那些王侯将相的坟，早变成了种庄稼的田。忙了一辈子鞠躬的人，到头来也不过一抔土。

及时行乐是古诗的老题目。《生年不满百》说"昼短苦夜长，何不秉烛游"，唐寅的花下眠，就是一千多年后的"秉烛游"。书里读过陶渊明的《归园田居》，也是归隐——但陶渊明是安静的，采菊东篱下，悠然见南山；唐寅是闹腾的，大张旗鼓地喝酒，到处跟人说我不干了。一个是淡，一个是赌气，赌气的那份也很真实。

我喜欢的正是他这股劲：被命运摔了一跤，没爬起来骂街，也没趴下不起，而是坐在桃花树下喝酒，把日子过成了一首诗。你学他一半就行——别人笑你的时候，你心里有自己的一朵桃花就行。酒就不用学了，你那杯奶茶就挺好。
\end{appreciation}
